\chapter{세상은 어떻게 안으로 들어오는가}
% full version: chapters/11_sensors.tex

모든 감각 기관은 마이크나 광전 소자 같은 변환기다. 빛이나 소리나 분자가 감각 세포 안의 특별한 단백질에 작용하면, 단백질이 “밸브”를 열고 전류가 흐른다. 이 사슬의 끝에서 딸깍이 만들어져 네트워크로 나간다. 감각 세포는 거의 모두 글루탐산을 내보내므로, 기계의 입력은 전화망에 곧바로 연결된다.

\section*{물리학의 한계에서}

뇌의 감지기는 가능한 것의 한계 바로 옆에서 일한다. 어둠 속에서 빛 감지기는 빛 알갱이 단 하나에도 반응한다. 소리 감지기는 1나노미터보다 작은, 원자 몇 개 크기보다 작은 변위를 알아챈다. 어스름 속에서는 정확도를 제한하는 것이 눈이 아니라 빛 자체다. 빛 알갱이는 무작위로 도착하므로 색조를 구별하려면 알갱이를 더 많이 모아야 한다. 그래서 해 질 녘에는 보는 속도가 느려진다. 노출 시간이 긴 카메라처럼 눈이 빛을 더 오래 모으기 때문이다.

\section*{자동 노출}

별이 뜬 밤에서 햇빛 비치는 눈밭까지 조도는 수십억 배 달라지고, 소리의 크기는 그보다도 더 크게 달라진다. 반면 딸깍의 빈도는 초당 0에서 수백 번 사이에서만 달라진다. 앞의 범위를 뒤의 범위에 그대로 밀어 넣을 수는 없다. 해법은 카메라와 같다. 자동 노출이다. 감지기는 주변의 평균 수준에 맞춰 감도를 끊임없이 조정하고, 밝기 자체가 아니라 \emph{배경에 비한} 밝기를 알린다. 시간이 지나도 변하지 않는 것에는 더 이상 반응하지 않고, 변화에는 매번 반응한다. 기계의 입력은 처음부터 변화에 대해 말한다.

\pencilfig{125}{%
% light rises in two tenfold steps, then drops a hundredfold; sensor rate = baseline + gain * (log light - adapted level),
% the adapted level follows log light with a 0.7 s time constant; rate clipped at zero
\draw[penlite=pcAmber] (0,1.75) -- (1.4,1.75) -- (1.4,2.1) -- (4.2,2.1) -- (4.2,2.45) -- (6.3,2.45) -- (6.3,1.75) -- (8.4,1.75);
\node[lbl, text=pcAmber!70!black, anchor=east] at (-0.05,2.0) {빛};
\node[lbl, anchor=south] at (2.8,2.12) {10배 밝게}; \node[lbl, anchor=south] at (5.25,2.47) {다시 10배};
\node[lbl, anchor=south] at (7.35,1.77) {다시 어둡게};
\draw[pen] (0,0) -- (8.4,0);
\draw[pcGray, densely dashed, line width=0.5pt] (0,0.32) -- (8.4,0.32);
\draw[penlite=pcBlue] plot coordinates {(0.00,0.32) (1.26,0.32) (1.40,1.02) (1.54,0.85) (1.68,0.72) (1.82,0.62) (1.96,0.54) (2.10,0.49) (2.24,0.45) (2.38,0.41) (2.52,0.39) (2.66,0.37) (2.80,0.36) (3.08,0.34) (3.50,0.33) (4.06,0.32) (4.20,1.03) (4.34,0.85) (4.48,0.72) (4.62,0.62) (4.76,0.54) (4.90,0.49) (5.04,0.45) (5.18,0.41) (5.32,0.39) (5.46,0.37) (5.60,0.36) (5.88,0.34) (6.16,0.33) (6.30,0.00) (7.00,0.00) (7.14,0.07) (7.28,0.13) (7.42,0.18) (7.56,0.22) (7.70,0.24) (7.84,0.26) (7.98,0.28) (8.12,0.29) (8.40,0.30)};
\node[lbl, text=pcBlue, anchor=east] at (-0.05,0.32) {응답};
\node[lbl, anchor=north] at (6.65,-0.02) {침묵};
}{빛이 한 단계에 열 배씩 계단처럼 밝아진다. 감지기는 변화가 있을 때마다 치솟았다가 금세 원래 수준으로 돌아온다. 다시 어두워지면 잠시 침묵한다.}

엔지니어들은 이 기법을 이벤트 카메라에 재현했다. 이벤트 카메라는 클록에 맞춰 프레임을 찍지 않는다. 무언가 바뀌면 각 픽셀이 스스로 “밝아졌다” 또는 “어두워졌다”를 알린다. 2장에 나온 “예”와 “아니요” 전선 한 쌍을 실리콘으로 만든 것이다.

\section*{눈은 그림을 압축한다}

눈에는 빛 감지기가 약 1억 개 있지만, 이미지를 뇌로 나르는 전선은 약 100만 개다. 그림은 눈을 떠나기 전에 약 100배로 압축된다. 균일한 영역에는 비용이 거의 들지 않고, 가장자리와 변화가 전달된다. 동물의 눈에서 얻은 값을 사람에 맞게 환산한 추정에 따르면, 눈은 뇌에 초당 약 1000만 비트를 보낸다. 오래된 네트워크 케이블 정도다.

\section*{귀는 피아노다}

귀는 엔지니어가 필터 뱅크를 놓듯이 소리를 주파수별로 나눈다. 귀 안의 탄성 있는 막은 자리마다 다른 주파수에 반응하고, 각 자리의 감지기는 자기 “건반”을 듣는다. 어느 감지기가 반응했느냐가 곧 음높이다.

\pencilfig{11}{\pencilart{11}}{귀는 거꾸로 된 피아노다. 건반마다 자기 음을 듣는다.}

게다가 귀는 수동적이지 않다. 귀의 세포 일부는 소리에 맞춰 스스로 막을 흔들어, 작은 소리는 수천 배 키우고 큰 소리는 거의 키우지 않는다. 자기 발진 직전의 경계에 놓인 증폭기이며, 바로 그곳에서 가장 민감하다. 낮은 주파수에서는 뉴런이 파동에 맞춰 딸깍하여 소리의 정확한 타이밍을 보존한다. 뇌는 이 타이밍으로 소리가 나는 방향을 찾는다. 주파수가 높아질수록 뉴런은 파동을 따라가지 못하고, 자리로 나타내는 부호만 남는다.

\section*{나머지 감각}

거의 모든 감각에서 앞 장들의 기법을 알아볼 수 있다. 온도는 따뜻함과 차가움에 따로 배정된 전선 두 개로 전달된다. 촉각은 닿는 순간 반응하고 곧 조용해지는 감지기와, 압력이 있는 동안 반응을 유지하는 감지기를 함께 쓴다. 후각은 바코드처럼 작동한다. 사람의 후각 감지기는 약 400가지 유형이 있고, 냄새는 반응한 유형들의 조합이다. 그런 조합의 수는 천문학적이다.

\fullversion{11}
